\begin{enunciado}{\ejExtra[10/06/26 🌶Ejercicio de la Práctica]}
  \fechaEjercicio{10/06/26 Ejercicio de la Práctica}
  Sean $a$, $b \en \enteros$ tal que $(a:b) = 7$.
  Hallar todos los $d = (15a^2 - 9b + 27: 189)$.
\end{enunciado}

Arranco \textit{coprimizando}, dado que tener números coprimos permite inferir la divisibilidad o no de las expresiones:
$$
  (a:b) = 7
  \Sii{
    $\llaves{rcl}{
        a = 7A\\
        b = 7B
      }$
  }
  d = 3 \cdot \ub{(5\cdot7^2A^2 - 3 \cdot 7 B + 9: 3^2 \cdot 7)}{D}
  \sii
  d = 3 \cdot D
$$
Estudio ahora los posibles valores de $D$ con $A \cop B$.

\bigskip
\parrafoDestacado{
  El \textit{máximo común divisor} se consigue haciendo el producto de las potencias,
  de la factorización en primos, comunes entre ambos números, de mayor exponte.
}

Posibles valores de $D$ dado que debe cumplir que:

$$
  \llave{rcl}{

    D & \divideA & 5 \cdot 7^2 A^2 - 3 \cdot 7 B + 9 \\
    & &\land  \\
    D & \divideA & 3^2 \cdot 7
  }
  \entonces
  D \en \set{1, 3, 7, 9, 21, 63}
$$
Estudio la divisibilidad de $ 5 \cdot 7^2 A^2 - 3 \cdot 7 B + 9$
para ver de \textit{descartar} potenciales valores de $D$.
Sale rápido que:
$$
  \llave{l}{
    \congruencia{ 5 \cdot 7^2 A^2 - 3 \cdot 7 B + 9}{2}{7} \\
    \congruencia{ 5 \cdot 7^2 A^2 - 3 \cdot 7 B + 9}{0}{3}
  }
  \Entonces{actualizo}[$D$]
  D \en \set{1, 3, 9}
$$
\medskip

\textit{Estudio divisibilidad por $9$:}
$$
  \llave{l}{
    \congruencia{ 5 \cdot 7^2 A^2 - 3 \cdot 7 B + 9}{2 \cdot (\ub{r_9(A)^2 + 3 \cdot r_9(B))}{\distinto 0 ~ ~ \text{para} ~~ A \cop B}}{9} \\
  }
$$
Para que la expresión sea nula, deberían tener tanto $A$ como $B$ a 3 como un divisor común y eso no puedo ocurrir dado que son \textit{coprimos}.

Actualizando $D$:
$$
  D \en \set{1,3}
$$

\medskip

\textit{Estudio divisibilidad por $3$:}

\magenta{Suponiendo el caso particular} en que $A$ \ul{no} es un múltiplo de 3:

$$
  \begin{array}{|r|c|c|c|} \hline
    r_3(A)                 & 1 & 2 \\\hline
    r_3((r_3(A)^2)         & 1 & 1 \\ \hline
    r_3(2 \cdot (r_3(A)^2) & 2 & 2 \\ \hline
  \end{array}
  \entonces
  \congruencia{5 \cdot 7^2 A^2 - 3 \cdot 7 B + 9}{\ub{2 \cdot (r_3(A))^2}{3 \noDivide A \entonces \distinto 0}}{3}
$$
Si no es divisible por 3, tampoco lo será por 9. Concluyo así que \textit{en este caso}:
$$
  D = 1 \entonces d = 3
$$

\bigskip

\magenta{Suponiendo el caso particular} en que $A$ \ul{sí} es un múltiplo de 3:
$$
  \text{Si} ~ ~ A = 3\alpha ~ \land ~ A \cop B \entonces \alpha \cop B \land 3 \noDivide B
$$
Puedo reescribir la expresión de $D$:
$$
  D = 3 \cdot \ub{(3 \cdot 5\cdot7^2 \alpha^2 - 7 B + 3: 3 \cdot 7)}{\delta} \sii D = 3 \delta
$$
\parrafoDestacado{
  \it
  ¿Qué sé sobre $\delta$?
}
Sin pensar nada sé que $\delta \en \set{1, 3, 7, 21}$, pero con un mínimo de ganas:
$$
  \begin{array}{lcl}
    \congruencia{3 \cdot 5 \cdot 7^2 \alpha^2 - 7 B + 3}{3}{7}                                & \entonces & \delta \distinto 7 \\
    \congruencia{3 \cdot 5 \cdot 7^2 \alpha^2 - 7 B + 3}{2 \cdot \ub{r_3(B)}{\distinto 0}}{3} & \entonces & \delta \distinto 3
  \end{array}
$$
Concluyo así que cuando $A = 3 \cdot \alpha$:
$$
  \delta = 1 \entonces D \divideA 3 \entonces d \divideA 9
$$

Por lo tanto:
$$
  \cajaResultado{
    d =
    \llave{rcl}{
      9 & \sii & 21 \divideA a \\
      3 & \multicolumn{2}{l}{\text{si no}}
    }
  }
$$

\begin{aportes}
  \item \aporte{\dirRepo}{naD GarRaz \github}
\end{aportes}
